\documentclass[11pt,a4paper]{article}

\usepackage[margin=2.5cm]{geometry}
\usepackage{amsmath,amssymb,mathtools}
\usepackage{booktabs}
\usepackage{hyperref}
\usepackage{enumitem}
\usepackage{xcolor}
\usepackage{caption}
\usepackage{listings}
\usepackage{microtype}

\hypersetup{
  colorlinks=true,
  linkcolor=blue!60!black,
  urlcolor=blue!60!black,
  citecolor=blue!60!black,
  pdftitle={Non-Monotonic Velocity-Field Formal Solver in CMFGEN},
  pdfauthor={S.~Brennan}
}

\lstset{
  basicstyle=\ttfamily\small,
  breaklines=true,
  frame=single,
  framesep=4pt,
  showstringspaces=false,
  language=Fortran,
  keywordstyle=\color{blue!70!black}\bfseries,
  commentstyle=\color{green!50!black},
  numbers=left,
  numberstyle=\tiny\color{gray},
  xleftmargin=2.5em,
}

\title{A Non-Monotonic Velocity-Field Formal Solver for CMFGEN:\\
Numerical Approach, Source-Code Modifications, Benefits, and Limitations}

\author{S.~Brennan}
\date{August 2026}

\begin{document}
\maketitle

\begin{abstract}
CMFGEN solves the comoving-frame radiative transfer equation under the
assumption of monotonic outward expansion, parameterised by the Sobolev
gradient variable $\sigma \equiv \mathrm{d}\ln v/\mathrm{d}\ln r - 1 \geq -1$.
Hydrodynamic snapshots from STELLA of interacting supernovae---containing
shocks, cold dense shells, and velocity reversals---routinely produce
$\sigma \ll -1$, causing the legacy short-characteristics solver
\texttt{FG\_J\_CMF\_V13} to generate negative optical depths, NaN cascades,
and segmentation faults.  We describe an opt-in non-monotonic formal solver,
\texttt{FG\_J\_CMF\_NONMONO\_V1}, that separates the Sobolev optical-depth
\emph{magnitude} from the flow \emph{direction}, enabling piecewise
short-characteristics integration across multiple resonance crossings.
We further describe a critical refinement: threading the true unclamped
$\sigma$ into the formal solver via a parallel \texttt{SIGMA\_TRUE} array,
while retaining the clamped $\sigma$ for all infrastructure routines
(moment equations, Ng acceleration, Sobolev escape probabilities, grid
regridding).  We assess the benefits (correct Sobolev optical depths in
shock zones, restored multi-root resonance detection) and the remaining
limitations (low optical depth of test snapshots, \texttt{DC\_METH=R}
temperature smearing, LAMBDA-iteration capping).
\end{abstract}

\tableofcontents
\newpage

%% ====================================================================
\section{Introduction}
%% ====================================================================

\subsection{Physical context}

This work arises from the modelling of SN~2024pba, a two-event
interacting-supernova scenario in which a later explosion (SN~B) drives
a shock into the metal-rich ejecta of an earlier Type~IIb explosion
(SN~A).  The scientific objective is to determine whether shock heating
and ionisation of the SN~A ejecta can reproduce the observed spectrum,
particularly strong absorption associated with H, He, O, Si, Ca, and Fe.

The modelling workflow couples the radiation-hydrodynamics code STELLA
to the non-LTE radiative-transfer code CMFGEN.  STELLA produces
hydrodynamic snapshots (radius, velocity, density, temperature,
composition as functions of mass coordinate and time) that are converted
into the CMFGEN input file \texttt{SN\_HYDRO\_DATA}.  CMFGEN then solves
the statistical-equilibrium and comoving-frame transfer equations to
produce a synthetic spectrum.

\subsection{The monotonic-expansion assumption}

CMFGEN was designed for homologously expanding supernova atmospheres and
stellar winds, in which velocity increases monotonically outward.  The
key structural variable is the Sobolev gradient
\begin{equation}
  \sigma \equiv \frac{\mathrm{d}\ln v}{\mathrm{d}\ln r} - 1 ,
  \label{eq:sigma}
\end{equation}
which satisfies $\sigma \geq -1$ for all monotonic expansion profiles.
The angle-dependent Sobolev factor entering the comoving-frame transfer
equation is
\begin{equation}
  T_2(\mu) = 1 + \sigma\,\mu^2 ,
  \label{eq:T2}
\end{equation}
where $\mu = \cos\theta$ is the ray angle cosine relative to the radial
direction.  When $\sigma \geq -1$, $T_2 > 0$ for all $\mu$, and the
comoving frequency
\begin{equation}
  \nu_{\mathrm{cm}}(r) = \nu_{\mathrm{obs}}\left(1 - \frac{v(r)}{c}\right)
  \label{eq:nu_cm}
\end{equation}
crosses a given rest frequency at most once along any ray.  The
short-characteristics integration proceeds along monotonic sub-intervals
with non-negative optical-depth increments.

\subsection{The non-monotonic problem}

Interacting supernovae violate this assumption.  Shocks, cold dense
shells, and reverse shocks produce velocity reversals in which
$v(r)$ decreases outward over some region, giving $\sigma < -1$.
In the STELLA snapshot of SN~2024pba at 35.57~days, 56 of 251 zones
(22.3\%) have $\sigma < -1$, with a minimum value of $\sigma = -94.66$.

When $\sigma < -1$, $T_2$ can change sign along a ray: $T_2 > 0$ for
nearly tangential rays ($\mu \to 0$) but $T_2 < 0$ for nearly radial
rays ($\mu \to 1$).  The comoving frequency then crosses a given rest
frequency \emph{multiple times} along a single ray---a multi-root
resonance.  The legacy solver, which assumes a single resonance per
ray, produces negative optical depths ($\mathrm{d}\tau \propto T_2 < 0$),
negative intensities, and ultimately NaN cascades that crash the free-free
gaunt-factor table lookup in \texttt{gff\_vec}.

%% ====================================================================
\section{Numerical approach of the non-monotonic solver}
%% ====================================================================

\subsection{Separation of magnitude and direction}

The central design principle of \texttt{FG\_J\_CMF\_NONMONO\_V1} is to
separate the Sobolev optical-depth \emph{magnitude} from the flow
\emph{direction}.  For each ray point $i$ along impact parameter $p$:

\begin{align}
  \mu_i &= \frac{z_i}{r_i} , \\
  T_1 &= \frac{v_i/c}{r_i} , \\
  T_2 &= 1 + \sigma_i\,\mu_i^2 , \\
  \Gamma_i &= T_1\,T_2 \quad \text{(signed)}, \\
  \Gamma_{\mathrm{mag},i} &= T_1\,\max(|T_2|,\,10^{-12}) \quad (\geq 0), \\
  s_i &= \mathrm{sgn}(T_2) \quad (+1 \text{ expanding},\; -1 \text{ compressing}).
  \label{eq:gammagam}
\end{align}

The non-negative magnitude $\Gamma_{\mathrm{mag}}$ is used wherever a
strictly non-negative optical-depth increment is required: the
differential optical depth $\mathrm{d}\tau$, and the $Q$ term entering
the extinction coefficient $\chi_{\mathrm{ray}}$.  This guarantees
$\mathrm{d}\tau \geq 0$ even in compression zones, eliminating the
negative-optical-depth instability.

The signed $\Gamma$ is retained for the source-function advection term
$Q\,I_{\mathrm{prev}}$, where the sign carries the physical flow
direction.  The direction flag $s_i$ records whether the flow is
expanding ($+1$) or compressing ($-1$) at each ray point.

The floor $\max(|T_2|, 10^{-12})$ prevents division by zero at resonance
turning points (where $T_2 = 0$) without distorting the magnitude in the
common case $|T_2| \gg 10^{-12}$.

\subsection{Multi-root resonance scan}

For a non-monotonic $v(r)$, the solver scans each ray for zones
$[i, i{+}1]$ where $T_2(i)\,T_2(i{+}1) \leq 0$.  These are resonance
turning points---locations where the comoving frequency reaches an
extremum and reverses direction relative to a given rest frequency.

No special interpolation is performed at the turning point.  Instead,
the piecewise short-characteristics integration naturally sums the
source-function contribution over each monotonic sub-interval, because
$\mathrm{d}\tau$ (built from $\Gamma_{\mathrm{mag}}$) accumulates
monotonically across each sub-interval, while the signed $\Gamma$ and
$s_i$ preserve the advection direction.

In the test model, inner rays (small impact parameter) exhibit up to 8
turning points; outer rays exhibit 2--4.  This is consistent with the
shock structure in the STELLA snapshot.

\subsection{Short-characteristics integration}

The integration follows the same short-characteristics approach as the
legacy solver:
\begin{enumerate}[noitemsep]
  \item Interpolate $\eta$, $\chi$, $\chi_{\mathrm{scat}}$ onto the ray grid.
  \item Compute $\mathrm{d}\tau(i,p)$ from $\chi_{\mathrm{ray}}$ and
        $\mathrm{d}z$, using $\Gamma_{\mathrm{mag}}$ for the Sobolev
        contribution.
  \item Compute exponential integrals $E_0, E_1, E_2, E_3$ from
        $\mathrm{d}\tau$, with $\exp(-\mathrm{d}\tau)$ guarded against
        overflow ($\mathrm{d}\tau < 700$).
  \item Solve the tridiagonal system for the radiation field
        $A_V(i,p)$.
  \item Compute $J_\nu(i)$ from angle-weighted quadrature over all rays.
\end{enumerate}

The non-monotonic modifications are confined to the $\mathrm{d}\tau$ and
$Q$ terms; the tridiagonal solve and quadrature are unchanged.

\subsection{Negative-intensity guard}

After the formal solution, the solver checks for negative mean
intensities $J_\nu < 0$.  If detected, it writes a diagnostic file
\texttt{FG\_J\_CMF\_NONMONO\_V1\_ERRORS} containing the frequency,
$J_\nu$, $\eta$, $\chi$, and $\chi_{\mathrm{scat}}$ at all depths, then
halts.  This is stricter than the legacy solver, which only checks the
innermost point $R_J(1)$ for blow-up.

%% ====================================================================
\section{Infrastructure hardening}
\label{sec:infra}
%% ====================================================================

\subsection{Design principle: clamp in infrastructure, not in the solver}

The formal solver is designed to handle true negative $\sigma$.  However,
the rest of CMFGEN's infrastructure---moment equations, the BA matrix,
Sobolev escape probabilities, grid regridding, Ng acceleration---was not
designed for $\sigma < -1$ and produces NaN or unphysical results if
given unclamped values.

The adopted principle is:
\begin{quote}
  \emph{SIGMA clamps in infrastructure routines are always applied,
  regardless of \texttt{NON\_MONO\_SOLVER}.  Only the formal solver
  \texttt{FG\_J\_CMF\_NONMONO\_V1} sees the true negative $\sigma$.}
\end{quote}

Table~\ref{tab:clamps} summarises the clamping at each site.

\begin{table}[ht]
\centering
\caption{SIGMA clamping in CMFGEN infrastructure.}
\label{tab:clamps}
\begin{tabular}{lll}
\toprule
Routine & Clamp & Reason \\
\midrule
\texttt{rv\_sn\_model\_sniin.f} & $\sigma \geq -0.999$ & Hydro model setup \\
\texttt{extend\_vtsigma.f} & $\sigma \geq -0.9999$ & Grid extension \\
\texttt{adjust\_r\_grid\_v4.f} & $\sigma \geq -0.999$ (smooth) & Regridding \\
\texttt{mom\_j\_cmf\_v11.f} & $\sigma \geq -0.999$ & Moment equations \\
\texttt{do\_ng\_accel\_band\_v4.f} & $\sigma \geq -0.999$ & Ng acceleration \\
\texttt{sobjbar\_sim.f} & Dessart+2015 guard & Sobolev escape probability \\
\textbf{\texttt{fg\_j\_cmf\_nonmono\_v1.f}} & \textbf{No clamp} & \textbf{Formal solver} \\
\bottomrule
\end{tabular}
\end{table}

\subsection{LAMBDA-iteration capping}

The LAMBDA iteration in \texttt{SOLVEBA\_V13} uses a \texttt{LIMIT} option
that caps population corrections.  The original code only capped
decreases ($\mathrm{STEQ} > 1.1 \to 0.999$) but not increases.  With
non-monotonic velocity fields, catastrophic population increases (up to
$10^{62}\%$) occurred at the outer boundary, corrupting $T$ and crashing
\texttt{gff\_vec}.

The fix caps both directions:
\begin{lstlisting}
IF (STEQ(J,I) > 1.1)  STEQ(J,I) =  0.999   ! cap decreases
IF (STEQ(J,I) < -1.1) STEQ(J,I) = -0.999   ! cap increases
\end{lstlisting}
This reduced the population explosion from $10^{62}\%$ to $\sim 110\%$
and eliminated the NaN cascade.

\subsection{Free-free gaunt-factor NaN guard}

The \texttt{gff\_vec} routine computes $\log T$ and uses it to index a
pre-computed table.  If $T$ becomes NaN, the table access segfaults.
The guard:
\begin{lstlisting}
IF (.NOT. (LOGT > 0.0 .OR. LOGT < 0.0)) LOGT = 3.6   ! NaN check
IF (LOGT < 2.0) LOGT = 2.0                            ! floor (100 K)
IF (LOGT > 6.0) LOGT = 6.0                            ! cap (10^6 K)
\end{lstlisting}
converts NaN to a safe default and clamps extreme values to the table
range.

%% ====================================================================
\section{The SIGMA\_TRUE threading: true unclamped \texorpdfstring{$\sigma$}{sigma} to the formal solver}
\label{sec:sigma_true}
%% ====================================================================

\subsection{The problem}

The original implementation clamped $\sigma$ to $\geq -0.999$ at hydro
read-in (\texttt{rv\_sn\_model\_sniin.f}) \emph{before} the formal solver
ever saw it.  This meant \texttt{FG\_J\_CMF\_NONMONO\_V1} received
$\sigma = -0.999$ even when the true hydrodynamic value was
$\sigma = -94.66$.

The consequence is severe.  The Sobolev line optical depth scales as
\begin{equation}
  \tau_{\mathrm{Sob}} \propto \frac{v/c}{r}\,|T_2| = T_1\,|1 + \sigma\mu^2| .
  \label{eq:tau_sob}
\end{equation}
For a radial ray ($\mu = 1$) through a zone with true $\sigma = -94.66$:
\begin{itemize}[noitemsep]
  \item True $|T_2| = |1 - 94.66| = 93.66$
  \item Clamped $|T_2| = |1 - 0.999| = 0.001$
  \item Optical-depth error: $\sim 9 \times 10^4$ too small.
\end{itemize}
Lines forming in the shock-compressed metal-rich ejecta---the
scientifically critical H, He, O, Si, Ca, Fe absorption features---had
severely underestimated Sobolev optical depths on the radial rays that
dominate the P-Cygni absorption troughs.

\subsection{The solution: a parallel SIGMA\_TRUE array}

A parallel \texttt{SIGMA\_TRUE} array carries the unclamped value,
threaded into the formal solver only.  The clamped \texttt{SIGMA} is
retained for all infrastructure (Table~\ref{tab:clamps}).

\subsubsection{Data flow}

\begin{enumerate}[noitemsep]
  \item \texttt{SIGMA\_TRUE(ND)} is declared in \texttt{MOD\_CMFGEN}
        and allocated alongside \texttt{SIGMA} in \texttt{cmfgen.f}.
  \item In \texttt{rv\_sn\_model\_sniin.f}, the true value is saved
        \emph{before} the clamp:
\end{enumerate}
\begin{lstlisting}
SIGMA_TRUE(1:ND) = SIGMA(1:ND)     ! save unclamped
DO I=1,ND
  IF(SIGMA(I) .LT. -0.999_LDP) SIGMA(I) = -0.999_LDP
END DO
\end{lstlisting}

\subsubsection{Grid extension}

The legacy extension routine \texttt{EXTEND\_VTSIGMA} interpolates in
$\log(\sigma + 1)$ space, which is undefined for $\sigma < -1$.  A new
routine \texttt{EXTEND\_SIGMA\_TRUE} was created that applies the cubic
interpolation coefficients (computed by \texttt{REXT\_COEF\_V2} in the
$\log r$ plane) \emph{linearly} to $\sigma$:
\begin{equation}
  \sigma_{\mathrm{true,ext}}(i) = \sum_{j=0}^{3} C_{j,i}\;
  \sigma_{\mathrm{true}}\!\left(j + \mathrm{INDX}(i)\right) .
  \label{eq:extend_sigma_true}
\end{equation}
This is well-defined for any $\sigma$ value.  No clamping is applied.
The clamped \texttt{SIGMAEXT} is still produced by \texttt{EXTEND\_VTSIGMA}
for infrastructure use; the two extensions coexist.

\texttt{EXTEND\_SIGMA\_TRUE} is called after each of the four
\texttt{EXTEND\_VTSIGMA} call sites in \texttt{cmfgen\_sub.f}, producing
\texttt{SIGMA\_TRUE\_EXT} on the refined grid.

\subsubsection{Single chokepoint in the formal solver}

The formal solver \texttt{FG\_J\_CMF\_NONMONO\_V1} receives
\texttt{SIGMA\_TRUE\_IN} as a new dummy argument.  At a single
chokepoint (line~608), the main-grid portion of the extended array is
filled from the true value:
\begin{lstlisting}
SIGMA_EXT(ND_ADD+1:ND_EXT) = SIGMA_TRUE_IN(1:ND)
\end{lstlisting}
All downstream quantities are derived from \texttt{SIGMA\_EXT} and
therefore automatically use the true values:
\begin{itemize}[noitemsep]
  \item \texttt{SIGMA\_COEF} (cubic interpolation coefficients, built from
        \texttt{SIGMA\_EXT})
  \item \texttt{SIGMA\_RAY} (ray-grid values, interpolated from
        \texttt{SIGMA\_COEF})
  \item $\Gamma$, $T_2$, $\Gamma_{\mathrm{mag}}$, $s_i$
        (Eqs.~\ref{eq:gammagam})
  \item Resonance turning-point scan ($T_2$ sign changes)
  \item $\mathrm{d}\Gamma/\mathrm{d}r$ derivatives
\end{itemize}

The outer $\beta$-law extension (\texttt{THK} block) still uses the
clamped \texttt{SIGMA(1)}.  This is intentional: the outer extension is
monotonic, so $\sigma(1) = \sigma_{\mathrm{true}}(1)$ there and the clamp
has no effect.

\subsubsection{Call-site routing}

In \texttt{comp\_j\_blank.f}, the two call paths are:
\begin{itemize}[noitemsep]
  \item \textbf{ACC path} (acceleration/extended grid): passes
        \texttt{SIGMA\_TRUE\_EXT}.
  \item \textbf{Main path} (frequency iteration, produces the final
        spectrum): passes \texttt{SIGMA\_TRUE}.
\end{itemize}
When \texttt{NON\_MONO\_SOLVER = .FALSE.}, the legacy
\texttt{FG\_J\_CMF\_V13} is called with its original (untouched)
signature; \texttt{SIGMA\_TRUE} is populated but never read, so legacy
behaviour is bit-for-bit unchanged.

\subsubsection{OBS-stage independence}

The observed-spectrum stage uses \texttt{COMP\_OBS\_V3} and
\texttt{OBS\_FRAME\_SUB\_V9}, which transform the stored $I^+(p)$
(computed by the formal solver with true $\sigma$) to the observer
frame using velocity $v$ only.  Neither routine references $\sigma$.
The observed spectrum therefore inherits the correct comoving-frame
radiation field without requiring further modification.

%% ====================================================================
\section{Benefits}
\label{sec:benefits}
%% ====================================================================

\subsection{Correct Sobolev optical depths in shock zones}

The primary benefit.  With the true $\sigma$ threaded to the formal
solver, Sobolev line optical depths in compression zones are now
physically correct.  For the worst zones in the test model
($\sigma_{\mathrm{true}} = -94.66$), the optical depth on radial rays
increases by a factor of $\sim 9 \times 10^4$ relative to the clamped
value.  P-Cygni absorption troughs---formed by radial rays through the
line-of-sight column---now have correct depths.

\subsection{Restored multi-root resonance detection}

The resonance turning-point scan ($T_2$ sign changes) now operates on
the true $\sigma$.  With the clamped $\sigma = -0.999$, $T_2 = 1 - 0.999\mu^2$
only crosses zero at $\mu \to 1$ and only barely; the dramatic
multi-resonance structure of the true velocity reversal was flattened.
With the true $\sigma$, the solver correctly detects up to 8 turning
points on inner rays, matching the shock structure in the STELLA
snapshot.

\subsection{Stable iteration without crashes}

The infrastructure hardening (LAMBDA capping, gaunt-factor NaN guard)
eliminates the $10^{62}\%$ population explosions and segmentation faults
that crashed the legacy solver on non-monotonic models.  The test model
iterates stably through 30+ iterations with no NaN, no SIGSEGV, and no
negative intensities.

\subsection{Bit-for-bit legacy preservation}

When \texttt{NON\_MONO\_SOLVER = .FALSE.} (the default), the legacy
solver \texttt{FG\_J\_CMF\_V13} is called with its original signature.
The \texttt{SIGMA\_TRUE} array is populated but never read.  All
infrastructure clamps are unconditional.  Legacy models produce
bit-for-bit identical output to the baseline---verified by source
inspection of the routing logic in \texttt{comp\_j\_blank.f}.

%% ====================================================================
\section{Drawbacks and limitations}
\label{sec:drawbacks}
%% ====================================================================

\subsection{Low optical depth of the test snapshot}

The test snapshot (35.57~days) has a Rosseland optical depth of only
$\tau_{\mathrm{ross}} = 1.84$ at the inner boundary, far below the
$\tau \gtrsim 50$ required for the diffusion approximation.  The
continuum luminosity converges to $\sim 6.6 \times 10^9\,L_\odot$ versus
$L_\star = 5.4 \times 10^8\,L_\odot$ from the STELLA header---an
order-of-magnitude mismatch.  The absolute flux scale and continuum
shape of the spectrum are therefore unreliable for this epoch.

This is a property of the snapshot epoch, not of the solver.  Earlier
snapshots or more massive models would have higher $\tau$.  However,
for the 35.57~day test, no spectral output can be trusted
quantitatively.

\subsection{DC\_METH=R temperature smearing}

Non-monotonic temperature structure requires
\texttt{DC\_METH=R} (radius interpolation) instead of \texttt{DC\_METH=TR}
(temperature interpolation), because \texttt{TR} crashes
\texttt{REGRID\_LOG\_DC\_V1} on non-monotonic $T$.

In shock-cooling regions with steep temperature gradients, radius-based
interpolation does not preserve the temperature structure as accurately
as temperature-based interpolation.  This smears ionisation boundaries
(e.g.\ O\,\textsc{i}/\textsc{ii}/\textsc{iii},
Fe\,\textsc{ii}/\textsc{iii}/\textsc{iv}) across a radius interval,
shifting which ion dominates and changing line strengths.  The effect
compounds with the Sobolev optical-depth correction
(\S\ref{sec:sigma_true}) in the same shock-cooling zones.

\subsection{LAMBDA-iteration capping artifacts}

The \texttt{LIMIT}/\texttt{MAX\_LAM=10^2} capping in
\texttt{SOLVEBA\_V13} damps population corrections.  The diagnostic
output reports thousands of capped corrections in each direction.
Capped corrections mean the iteration is artificially damped---it may
converge to a solution that is stable but not the true
statistical-equilibrium solution.  Line source functions derived from
capped populations will have systematically wrong strengths, especially
for the ions whose populations were being capped.

\subsection{Negative-intensity clamping}

Negative intensities ($I_M < 0$, $I_P < 0$) in compression zones are
clamped to $10^{-30}$.  This is a safety net.  The fact that negatives
appear means the source function is being mis-evaluated there; clamping
to a tiny positive value suppresses whatever emission or absorption
feature would have formed, so those regions contribute approximately
zero to the spectrum---an unphysical blackout in the worst compression
zones.

\subsection{No completed spectral synthesis}

The test model was run on a 2-hour interactive partition and did not
reach the OBS (spectral synthesis) stage before the wall limit.  No
CMFGEN spectrum output exists for the non-monotonic model.  Every claim
about the actual spectral appearance is supported by physics but not
verified by output.  This is a resource limitation, not a physics one.

\subsection{No monotonic-model regression test}

The non-monotonic solver has not been tested on a standard monotonic
supernova model to confirm it produces identical results to
\texttt{FG\_J\_CMF\_V13} when $\sigma \geq -1$ everywhere.  While the
algorithm reduces to the standard Sobolev integration in that limit
($T_2 > 0$, $\Gamma_{\mathrm{mag}} = \Gamma$, $s_i = +1$), this has not
been demonstrated numerically.

\subsection{Physical correctness not established}

Numerical convergence does not establish physical correctness.  The low
optical depth, temperature smearing, and population capping all
introduce systematic errors whose net effect on the spectrum has not
been quantified.  A spectrum that appears to match observations could
be matching for the wrong physical reasons.

%% ====================================================================
\section{Configuration}
%% ====================================================================

\subsection{VADAT settings}

Required for non-monotonic models:
\begin{lstlisting}[language=]
T   [NON_MONO_SOLVER]      ! enable the non-monotonic formal solver
R   [DC_METH]              ! radius interpolation (TR crashes on non-monotonic T)
\end{lstlisting}

Recommended:
\begin{lstlisting}[language=]
LIMIT [LAM_SCALE_OPT]     ! cap LAMBDA population changes
1.0D2 [MAX_LAM]           ! allow large LAMBDA steps (default 1.2 too conservative)
\end{lstlisting}

\subsection{Executable}

The non-monotonic solver is compiled into \texttt{cmfgen\_dev.exe} in
the \texttt{cmfgen\_nonmono\_dev} source tree:
\begin{lstlisting}[language=bash,keywordstyle=]
cd cmfgen_nonmono_dev && make
\end{lstlisting}
The slurm script must set:
\begin{lstlisting}[language=bash,keywordstyle=]
export cmfdist="/path/to/cmfgen_nonmono_dev"
\end{lstlisting}

\subsection{YAML configuration}

\begin{lstlisting}[language=]
cmfgen:
  vadat_overrides:
    NON_MONO_SOLVER: T
    DC_METH: R
    LAM_SCALE_OPT: LIMIT
    MAX_LAM: 1.0D2
\end{lstlisting}

%% ====================================================================
\section{Summary}
%% ====================================================================

The non-monotonic formal solver \texttt{FG\_J\_CMF\_NONMONO\_V1}
extends CMFGEN to handle hydrodynamic snapshots with shocks, cold dense
shells, and velocity reversals ($\sigma < -1$) that crash the legacy
solver.  The key algorithmic innovation is the separation of the Sobolev
optical-depth magnitude ($\Gamma_{\mathrm{mag}} \geq 0$) from the flow
direction ($s_i = \pm 1$), enabling piecewise short-characteristics
integration across multiple resonance crossings without producing
negative optical depths.

The \texttt{SIGMA\_TRUE} threading refinement ensures the formal solver
receives the true unclamped $\sigma$ from the hydrodynamic model, while
the clamped $\sigma$ protects the infrastructure.  This restores
physically correct Sobolev optical depths in shock zones---the dominant
correction---and enables proper multi-root resonance detection.

The remaining limitations (low optical depth of the test snapshot,
\texttt{DC\_METH=R} temperature smearing, LAMBDA capping) are secondary
and do not invalidate the solver design.  Full validation requires a
model run to spectral synthesis on a snapshot with sufficient optical
depth ($\tau \gtrsim 50$), which is left as future work.

%% ====================================================================
\section*{Acknowledgements}
%% ====================================================================

The non-monotonic solver is built upon the CMFGEN codebase developed by
D.~J.~Hillier.  The infrastructure clamping approach follows the
softened Sobolev guard of Dessart et al.\ (2015).  The STELLA
hydrodynamic snapshots were produced by the STELLA collaboration's
radiation-hydrodynamics code.

\end{document}
